% Options for packages loaded elsewhere
\PassOptionsToPackage{unicode}{hyperref}
\PassOptionsToPackage{hyphens}{url}
\PassOptionsToPackage{dvipsnames,svgnames,x11names}{xcolor}
%
\documentclass[
]{article}
\usepackage{amsmath,amssymb}
\usepackage{iftex}
\ifPDFTeX
  \usepackage[T1]{fontenc}
  \usepackage[utf8]{inputenc}
  \usepackage{textcomp} % provide euro and other symbols
\else % if luatex or xetex
  \usepackage{unicode-math} % this also loads fontspec
  \defaultfontfeatures{Scale=MatchLowercase}
  \defaultfontfeatures[\rmfamily]{Ligatures=TeX,Scale=1}
\fi
\usepackage{lmodern}
\ifPDFTeX\else
  % xetex/luatex font selection
\fi
% Use upquote if available, for straight quotes in verbatim environments
\IfFileExists{upquote.sty}{\usepackage{upquote}}{}
\IfFileExists{microtype.sty}{% use microtype if available
  \usepackage[]{microtype}
  \UseMicrotypeSet[protrusion]{basicmath} % disable protrusion for tt fonts
}{}
\makeatletter
\@ifundefined{KOMAClassName}{% if non-KOMA class
  \IfFileExists{parskip.sty}{%
    \usepackage{parskip}
  }{% else
    \setlength{\parindent}{0pt}
    \setlength{\parskip}{6pt plus 2pt minus 1pt}}
}{% if KOMA class
  \KOMAoptions{parskip=half}}
\makeatother
\usepackage{xcolor}
\usepackage[margin=1in]{geometry}
\usepackage{color}
\usepackage{fancyvrb}
\newcommand{\VerbBar}{|}
\newcommand{\VERB}{\Verb[commandchars=\\\{\}]}
\DefineVerbatimEnvironment{Highlighting}{Verbatim}{commandchars=\\\{\}}
% Add ',fontsize=\small' for more characters per line
\newenvironment{Shaded}{}{}
\newcommand{\AlertTok}[1]{\textcolor[rgb]{1.00,0.00,0.00}{\textbf{#1}}}
\newcommand{\AnnotationTok}[1]{\textcolor[rgb]{0.38,0.63,0.69}{\textbf{\textit{#1}}}}
\newcommand{\AttributeTok}[1]{\textcolor[rgb]{0.49,0.56,0.16}{#1}}
\newcommand{\BaseNTok}[1]{\textcolor[rgb]{0.25,0.63,0.44}{#1}}
\newcommand{\BuiltInTok}[1]{\textcolor[rgb]{0.00,0.50,0.00}{#1}}
\newcommand{\CharTok}[1]{\textcolor[rgb]{0.25,0.44,0.63}{#1}}
\newcommand{\CommentTok}[1]{\textcolor[rgb]{0.38,0.63,0.69}{\textit{#1}}}
\newcommand{\CommentVarTok}[1]{\textcolor[rgb]{0.38,0.63,0.69}{\textbf{\textit{#1}}}}
\newcommand{\ConstantTok}[1]{\textcolor[rgb]{0.53,0.00,0.00}{#1}}
\newcommand{\ControlFlowTok}[1]{\textcolor[rgb]{0.00,0.44,0.13}{\textbf{#1}}}
\newcommand{\DataTypeTok}[1]{\textcolor[rgb]{0.56,0.13,0.00}{#1}}
\newcommand{\DecValTok}[1]{\textcolor[rgb]{0.25,0.63,0.44}{#1}}
\newcommand{\DocumentationTok}[1]{\textcolor[rgb]{0.73,0.13,0.13}{\textit{#1}}}
\newcommand{\ErrorTok}[1]{\textcolor[rgb]{1.00,0.00,0.00}{\textbf{#1}}}
\newcommand{\ExtensionTok}[1]{#1}
\newcommand{\FloatTok}[1]{\textcolor[rgb]{0.25,0.63,0.44}{#1}}
\newcommand{\FunctionTok}[1]{\textcolor[rgb]{0.02,0.16,0.49}{#1}}
\newcommand{\ImportTok}[1]{\textcolor[rgb]{0.00,0.50,0.00}{\textbf{#1}}}
\newcommand{\InformationTok}[1]{\textcolor[rgb]{0.38,0.63,0.69}{\textbf{\textit{#1}}}}
\newcommand{\KeywordTok}[1]{\textcolor[rgb]{0.00,0.44,0.13}{\textbf{#1}}}
\newcommand{\NormalTok}[1]{#1}
\newcommand{\OperatorTok}[1]{\textcolor[rgb]{0.40,0.40,0.40}{#1}}
\newcommand{\OtherTok}[1]{\textcolor[rgb]{0.00,0.44,0.13}{#1}}
\newcommand{\PreprocessorTok}[1]{\textcolor[rgb]{0.74,0.48,0.00}{#1}}
\newcommand{\RegionMarkerTok}[1]{#1}
\newcommand{\SpecialCharTok}[1]{\textcolor[rgb]{0.25,0.44,0.63}{#1}}
\newcommand{\SpecialStringTok}[1]{\textcolor[rgb]{0.73,0.40,0.53}{#1}}
\newcommand{\StringTok}[1]{\textcolor[rgb]{0.25,0.44,0.63}{#1}}
\newcommand{\VariableTok}[1]{\textcolor[rgb]{0.10,0.09,0.49}{#1}}
\newcommand{\VerbatimStringTok}[1]{\textcolor[rgb]{0.25,0.44,0.63}{#1}}
\newcommand{\WarningTok}[1]{\textcolor[rgb]{0.38,0.63,0.69}{\textbf{\textit{#1}}}}
\setlength{\emergencystretch}{3em} % prevent overfull lines
\providecommand{\tightlist}{%
  \setlength{\itemsep}{0pt}\setlength{\parskip}{0pt}}
\setcounter{secnumdepth}{-\maxdimen} % remove section numbering
\setlength{\emergencystretch}{3em}
\ifLuaTeX
  \usepackage{selnolig}  % disable illegal ligatures
\fi
\usepackage{bookmark}
\IfFileExists{xurl.sty}{\usepackage{xurl}}{} % add URL line breaks if available
\urlstyle{same}
\hypersetup{
  colorlinks=true,
  linkcolor={black},
  filecolor={Maroon},
  citecolor={Blue},
  urlcolor={black},
  pdfcreator={LaTeX via pandoc}}

\author{}
\date{}

\begin{document}

\section{BA1-T1 - Minimal BAE ontology candidate
derivation}\label{ba1-t1---minimal-bae-ontology-candidate-derivation}

\textbf{Revision:} R1

\textbf{Status:} COMPLETED / PROVISIONAL CANDIDATE / BA1 NOT CLOSED

\textbf{Repository baseline reviewed:}
\texttt{83af68cb1a02a6b1e76f591d4c1519f9496be3b3}

\textbf{Phase:} BA1 - Minimal BAE ontology

\textbf{BA0:} CLOSED

\textbf{BA2:} NOT STARTED

\subsection{1. Question and guardrail}\label{1-question-and-guardrail}

BA1-T1 asks only:

\begin{quote}
What is the smallest falsifiable set of first-class analytical identity
families needed to realize the closed BA0 responsibilities without
importing document types, tool structures, threat-method categories or
general systems-modeling taxonomies by default?
\end{quote}

The governing rule is:

\begin{Shaded}
\begin{Highlighting}[]
\NormalTok{semantic responsibility exists}
\NormalTok{        !=}
\NormalTok{first{-}class metaclass required}
\end{Highlighting}
\end{Shaded}

For every recurring semantic responsibility, BA1-T1 tests whether it can
remain a relation, role, property/value, constraint/state, provenance or
diagnostic record, or derived projection before proposing a first-class
BAE identity.

This microstep does not define relation vocabulary, canonical action
predicates, identity/equivalence lifecycle rules, provenance
materialization, projection contracts, AnalysisRecord, Finding or a
threat-method overlay.

\subsection{2. Closed input boundary from
BA0}\label{2-closed-input-boundary-from-ba0}

BA0-T3 closed eight responsibilities:

\begin{itemize}
\tightlist
\item
  \texttt{BA0-C1} authority and provenance boundary;
\item
  \texttt{BA0-C2} baseline-scoped shared semantic identity;
\item
  \texttt{BA0-C3} method-neutral shared core;
\item
  \texttt{BA0-C4} explicit uncertainty and diagnostic localization;
\item
  \texttt{BA0-C5} projection readiness and source drill-down;
\item
  \texttt{BA0-C6} change-impact traceability;
\item
  \texttt{BA0-C7} source-localized feedback handoff;
\item
  \texttt{BA0-C8} minimality and representation independence.
\end{itemize}

BA0 also closed that grounded, derived and diagnostic/unresolved are
required \textbf{origin/state responsibilities}, not pre-approved domain
metaclasses. A generic reviewed analytical-addition class is not
required.

The BA0 boundary already speaks of stable analytical \textbf{referents
and propositions}. BA1-T1 treats those words as responsibility-level
evidence and tests whether they deserve distinct first-class identity
families.

\subsection{3. Evidence base}\label{3-evidence-base}

\subsubsection{3.1 Facial-access branch inherited from
BA0-T1/T2}\label{31-facial-access-branch-inherited-from-ba0-t1t2}

The branch requires stable recognition of at least:

\begin{itemize}
\tightlist
\item
  \texttt{CameraSubsystem};
\item
  \texttt{RecognitionProcessor};
\item
  \texttt{RecognitionCapture};
\item
  \texttt{RecognitionRequest};
\item
  consumed transport responsibility/context;
\item
  current Ethernet realization;
\item
  the delivery/correlation/failure statements around \texttt{FR-3.4};
\item
  the confidentiality, integrity and authorized-origin conditions
  constraining that delivery.
\end{itemize}

BA0-T1 showed that the same referent can participate in grounded and
derived assertions with different origin. It also showed that a
plausible invented endpoint is not accepted merely because it is
convenient to draw.

Mutation M1 (\texttt{Ethernet\ -\textgreater{}\ Wi-Fi}) changes one
realization while preserving most other meaning. Mutation M3 (remote
-\textgreater{} local recognition) retires/revises delivery meaning and
proves that lexical equality alone cannot establish analytical identity.

\subsubsection{3.2 Independent order-fulfillment
corpus}\label{32-independent-order-fulfillment-corpus}

The revision-3 complete order-fulfillment authoring probe is used as a
structurally distinct pressure source. It explicitly states that its
governed references are an authoring substrate and that their final Base
Analysis types are deferred.

The corpus contains:

\begin{itemize}
\tightlist
\item
  human/external participants (\texttt{Customer},
  \texttt{WarehouseOperator}, \texttt{PaymentProvider});
\item
  project capabilities/services such as \texttt{OrderService},
  \texttt{InventoryService} and \texttt{PaymentAdapter};
\item
  further service referents such as \texttt{FulfillmentService} and
  \texttt{CarrierAdapter};
\item
  lifecycle-bearing concepts such as \texttt{OrderEvaluation},
  \texttt{Order} and \texttt{Reservation};
\item
  result concepts such as \texttt{ReservationResult} and
  \texttt{PaymentAuthorizationResult};
\item
  task concepts such as \texttt{FulfillmentTask};
\item
  store candidates such as \texttt{InventoryLedger} and
  \texttt{ReservationStore};
\item
  contract candidates such as \texttt{PaymentProviderContract} and
  \texttt{CarrierContract};
\item
  state and result values (\texttt{ACCEPTED}, \texttt{authorized},
  \texttt{declined}, \texttt{indeterminate}, reservation lifecycle
  state);
\item
  cross-branch service consumption without ownership transfer;
\item
  decisions that change responsibility boundaries (internal inventory
  authority versus external WMS);
\item
  operational statements about correlation, atomicity, idempotency,
  compensation, concurrency, failure and physical handoff milestones.
\end{itemize}

The corpus deliberately leaves the exact formal object for cross-MR
service consumption open for Base Analysis.

\subsubsection{3.3 Documentation-layer
constraints}\label{33-documentation-layer-constraints}

The FunctionalRequirement baseline makes normative prose semantically
primary and uses SPO references only to expose reusable participants. It
explicitly states that SPO alone does not contain all conditional,
concurrency, lifecycle or failure semantics.

The Decision layer records project commitments; it is not automatically
an analytical type. The MR working evidence likewise warns that nouns in
MR prose do not directly create BAEs. The layering work separates L1
semantics from material representation, tooling and projections.

Therefore BA1 must not copy the
\texttt{MR\ -\textgreater{}\ Decision\ -\textgreater{}\ FR\ -\textgreater{}\ SR}
document hierarchy into Base Analysis.

\subsection{4. Lower-bound falsification: can one generic element type
be
enough?}\label{4-lower-bound-falsification-can-one-generic-element-type-be-enough}

\subsubsection{Candidate alternative A}\label{candidate-alternative-a}

\begin{Shaded}
\begin{Highlighting}[]
\NormalTok{BAElement}
\end{Highlighting}
\end{Shaded}

with every project thing and every analytical claim represented as the
same undifferentiated kind.

\subsubsection{Test}\label{test}

This collapses two responsibilities that current evidence needs to
distinguish:

\begin{enumerate}
\def\labelenumi{\arabic{enumi}.}
\tightlist
\item
  \textbf{identity of what the project statement is about}; and
\item
  \textbf{identity, origin and evolution of the statement itself}.
\end{enumerate}

In T1, \texttt{CameraSubsystem} can remain the same referent while the
statements involving it are retained, revised or retired. The normalized
delivery tuple is derived even though its referenced participants are
grounded. In the order corpus, the same \texttt{OrderEvaluation}
identity participates in multiple independently governed claims about
reservation, payment, compensation and commitment.

If referent and proposition are one undifferentiated kind, provenance
and change review can still be implemented with extra flags, but the
semantic distinction has merely been hidden in properties. That does not
reduce the required ontology; it obscures it.

\subsubsection{Disposition}\label{disposition}

\textbf{REJECTED AS THE MINIMAL BA1-T1 CANDIDATE.} A single
undifferentiated analytical element does not preserve the observed
distinction between identity of a referent and identity/origin of claims
about it.

\subsection{5. Upper-bound falsification: are domain-specific roots
already
forced?}\label{5-upper-bound-falsification-are-domain-specific-roots-already-forced}

The reviewed evidence repeatedly uses structural participants, behavior,
information, state, interfaces/contracts, stores and boundaries. BA0-R
also established their recurring semantic relevance.

However, recurring relevance is not sufficient for first-class split.

Current cases can be represented without yet declaring separate root
metaclasses:

\begin{itemize}
\tightlist
\item
  \texttt{CameraSubsystem}, \texttt{RecognitionProcessor},
  \texttt{OrderService}, \texttt{PaymentProvider},
  \texttt{InventoryLedger} and \texttt{PaymentProviderContract} can all
  be independently identified referents whose more specific semantic
  role remains explicit but not yet a metaclass split;
\item
  \texttt{RecognitionCapture}, \texttt{RecognitionRequest},
  \texttt{OrderEvaluation}, \texttt{Reservation} and result objects can
  be referents with information/resource/lifecycle roles;
\item
  delivery, correlation, dependency, ownership, state transition,
  atomicity, idempotency, applicability and security conditions can be
  propositions;
\item
  a behavior/event that later needs independent cross-statement identity
  may itself become a referent, without proving today that all behaviors
  require a separate \texttt{Behavior} metaclass;
\item
  state values can remain values/qualifiers or propositions until an
  independently governed state identity is forced;
\item
  a boundary, channel, store or contract can remain a referent role/kind
  until evidence requires subtype-specific invariants that cannot be
  expressed without a first-class split.
\end{itemize}

\subsubsection{Disposition}\label{disposition-1}

\textbf{NOT YET FORCED.} \texttt{Behavior}, \texttt{Information},
\texttt{Participant}, \texttt{Component/Capability}, \texttt{Boundary},
\texttt{Channel}, \texttt{Store}, \texttt{Contract},
\texttt{State/Mode/Context}, \texttt{Interface/Connection},
\texttt{Dependency} and similar families remain split hypotheses for
later pressure testing.

\subsection{6. Minimal ontology
candidate}\label{6-minimal-ontology-candidate}

BA1-T1 therefore proposes exactly two first-class analytical identity
families.

\subsubsection{\texorpdfstring{BA1-C1 -
\texttt{BAReferent}}{BA1-C1 - BAReferent}}\label{ba1-c1---bareferent}

\textbf{Status:} CANDIDATE / FIRST-CLASS IDENTITY PRESSURE SUPPORTED /
NOT ACCEPTED

A \texttt{BAReferent} is an independently identifiable unit of
methodology-neutral shared project meaning that must be recognizable
across one or more analytical propositions, projections or governed
baselines.

A referent may denote structural, human/external,
informational/resource, behavioral/event, contractual, storage or other
project meaning. Those semantic kinds are not first-class subtype
decisions in BA1-T1.

Examples from the evidence include, when independent identity is needed
by analysis:

\begin{itemize}
\tightlist
\item
  \texttt{CameraSubsystem}, \texttt{RecognitionProcessor};
\item
  \texttt{RecognitionCapture}, \texttt{RecognitionRequest};
\item
  \texttt{OrderEvaluation}, \texttt{Reservation};
\item
  \texttt{InventoryService}, \texttt{PaymentProvider};
\item
  \texttt{PaymentProviderContract}, \texttt{InventoryLedger}.
\end{itemize}

A governed document is not automatically a \texttt{BAReferent}. The
project meaning expressed by a document may justify one or more
referents, while the document remains source authority and provenance
evidence.

\subsubsection{\texorpdfstring{BA1-C2 -
\texttt{BAProposition}}{BA1-C2 - BAProposition}}\label{ba1-c2---baproposition}

\textbf{Status:} CANDIDATE / FIRST-CLASS IDENTITY PRESSURE SUPPORTED /
NOT ACCEPTED

A \texttt{BAProposition} is an independently identifiable
methodology-neutral analytical assertion expressing shared project
meaning about one or more referents, with enough identity to carry
source/origin distinction, be retained/revised/retired or diagnosed, and
be reused by multiple projections or consumers.

A proposition may express, without pre-closing the later vocabulary:

\begin{itemize}
\tightlist
\item
  relation or dependency;
\item
  behavior/action/transfer/transition;
\item
  responsibility or ownership context;
\item
  correlation;
\item
  state/condition/result semantics;
\item
  applicability;
\item
  constraint/property;
\item
  lifecycle/failure/concurrency semantics.
\end{itemize}

A proposition may itself need to be the target of another proposition.
For example, a specialized condition may constrain the same delivery
meaning without forcing a separate \texttt{Behavior} metaclass at this
stage.

The exact predicate/action vocabulary and structural shape of
propositions belong to BA2. Exact origin/provenance/source-locator
materialization belongs to BA3.

\subsubsection{Candidate picture}\label{candidate-picture}

\begin{Shaded}
\begin{Highlighting}[]
\NormalTok{Base Analysis for one governed baseline}

\NormalTok{    BAReferent *}
\NormalTok{        \^{}}
\NormalTok{        | referenced by}
\NormalTok{        |}
\NormalTok{    BAProposition *}

\NormalTok{BAProposition has identity and may also be targeted by other propositions.}
\end{Highlighting}
\end{Shaded}

\texttt{BAE} is used as the umbrella term for Base Analysis elements.
BA1-T1 does \textbf{not} require an additional concrete or abstract
\texttt{BAElement} metaclass unless a later shared-field/invariant
analysis justifies it.

\subsection{7. Why proposition identity is not just a relation
table}\label{7-why-proposition-identity-is-not-just-a-relation-table}

The evidence requires more than binary edges.

\texttt{FR-3.4} includes source, destination, delivered information,
request correlation and failure semantics. The order corpus includes
conditions, atomicity, concurrency, idempotency, lifecycle and
indeterminate outcomes that cannot be reduced faithfully to one
subject-predicate-object edge.

SPO remains useful as one structured projection, but the full analytical
claim is the semantic unit that needs provenance and change review.

Therefore BA1-T1 does not define \texttt{Relation} as a first-class BAE.
A relation/action may be part of the content of a
\texttt{BAProposition}; BA2 will decide the minimal relation and action
vocabulary.

\subsection{8. Mapping the BA0 responsibilities to the
candidate}\label{8-mapping-the-ba0-responsibilities-to-the-candidate}

\begin{itemize}
\tightlist
\item
  \textbf{C1 authority/provenance:} referents and especially
  propositions remain source-aligned; exact provenance records are BA3,
  not new domain BAEs.
\item
  \textbf{C2 shared semantic identity:} \texttt{BAReferent} and
  \texttt{BAProposition} provide the two identity-bearing families
  forced by current evidence.
\item
  \textbf{C3 method-neutral core:} neither family contains
  STRIDE/STRIDE-AI or tool-owned categories.
\item
  \textbf{C4 diagnostic localization:} diagnostic/unresolved remains
  status/record over affected identities; no \texttt{Diagnostic} domain
  BAE is required.
\item
  \textbf{C5 projection readiness:} views select/project referents and
  propositions; \texttt{View} is BA4, not a core BAE type.
\item
  \textbf{C6 change-impact traceability:} referent/proposition identity
  allows selective retain/revise/retire review; exact dependency
  relations are BA2/BA3.
\item
  \textbf{C7 feedback handoff:} source-aligned propositions localize the
  governed area to review; governance workflow remains outside BA.
\item
  \textbf{C8 minimality/independence:} no notation, storage technology,
  document hierarchy or domain-specific subtype is accepted without
  separate pressure.
\end{itemize}

\subsection{9. Explicit non-first-class dispositions in
BA1-T1}\label{9-explicit-non-first-class-dispositions-in-ba1-t1}

\subsubsection{Not copied from
documentation}\label{not-copied-from-documentation}

The following are \textbf{not BAE metaclasses by inheritance}:

\begin{itemize}
\tightlist
\item
  \texttt{MacroRequirement};
\item
  \texttt{Decision};
\item
  \texttt{Requirement};
\item
  \texttt{FunctionalRequirement};
\item
  \texttt{SpecializedRequirement};
\item
  \texttt{SecurityRequirement};
\item
  \texttt{NormativeClause}.
\end{itemize}

They remain governed-document concepts/source identities. BA3 will later
define document-to-Base-Analysis derivation and provenance.

\subsubsection{Not promoted from historical/tool
studies}\label{not-promoted-from-historicaltool-studies}

The following are \textbf{not accepted roots}:

\begin{itemize}
\tightlist
\item
  \texttt{actor};
\item
  \texttt{component};
\item
  \texttt{asset};
\item
  \texttt{boundary};
\item
  \texttt{data\_flow};
\item
  ThreatForge-native objects;
\item
  STRIDE DFD categories.
\end{itemize}

\subsubsection{Deferred semantic-kind
splits}\label{deferred-semantic-kind-splits}

The following remain \textbf{OPEN / DEFERRED AS FIRST-CLASS SPLITS}:

\begin{itemize}
\tightlist
\item
  Participant / actor-like role;
\item
  structural element / component / capability;
\item
  information / resource;
\item
  behavior / event / transition;
\item
  interface / connection / channel;
\item
  store / persistence;
\item
  contract;
\item
  boundary / scope boundary;
\item
  state / mode / context;
\item
  dependency / allocation;
\item
  property / constraint as an autonomous element.
\end{itemize}

They may remain classifications/roles/propositions, or later evidence
may force a subtype/metaclass split.

\subsubsection{Deferred to later phases rather than domain
ontology}\label{deferred-to-later-phases-rather-than-domain-ontology}

\begin{itemize}
\tightlist
\item
  source locator, provenance record and origin materialization
  -\textgreater{} BA3;
\item
  diagnostic record/status schema -\textgreater{} BA3;
\item
  view/projection -\textgreater{} BA4;
\item
  controlled vocabulary and lexical assistance -\textgreater{} BA5;
\item
  AnalysisRecord, Finding and methodology output -\textgreater{}
  post-BA6 analysis envelope.
\end{itemize}

\subsection{10. Candidate falsification
criteria}\label{10-candidate-falsification-criteria}

The two-family candidate must be weakened or split if a concrete corpus
or consumer demonstrates any of the following:

\begin{enumerate}
\def\labelenumi{\arabic{enumi}.}
\tightlist
\item
  two referents with the same generic representation require
  subtype-specific invariants that cannot be expressed without losing
  reusable semantics;
\item
  behavior/event identity cannot be represented honestly as referent
  plus propositions without ambiguity or unstable reification;
\item
  information/resource identity requires distinct lifecycle/typing rules
  that generic referents cannot preserve;
\item
  state/context must be independently identified and related across
  multiple claims in a way that values/qualifiers/propositions cannot
  represent cleanly;
\item
  boundary/interface/contract/store semantics require independent
  first-class identity and subtype invariants for multiple analysis
  consumers;
\item
  proposition-as-target becomes ambiguous or creates circular semantics
  that a dedicated behavior/constraint element would resolve materially;
\item
  a method-neutral consumer cannot build its required projection without
  reconstructing domain kinds independently from raw prose;
\item
  a structurally different governed corpus cannot be represented without
  adding method/tool/document-specific meaning to either family.
\end{enumerate}

Ordinary requests for a nicer diagram or convenient implementation class
are not sufficient evidence.

\subsection{11. Result}\label{11-result}

BA1-T1 produces a \textbf{minimal falsifiable ontology candidate}, not
an accepted ontology:

\begin{Shaded}
\begin{Highlighting}[]
\NormalTok{CANDIDATE}
\NormalTok{    BAReferent}
\NormalTok{    BAProposition}

\NormalTok{NOT ACCEPTED YET}
\NormalTok{    any domain{-}specific subtype/metaclass split}
\end{Highlighting}
\end{Shaded}

This candidate is smaller than the historical five-focus study and
smaller than a general-purpose systems model. It is also stronger than
an untyped node graph because it preserves the distinction between
project referent identity and provenance-bearing analytical claim
identity.

No relation vocabulary is closed by BA1-T1.

\subsection{12. Next authorized
microstep}\label{12-next-authorized-microstep}

Execute only:

\begin{quote}
\textbf{BA1-T2 - split-or-collapse pressure test of the two-family
candidate.}
\end{quote}

BA1-T2 must deliberately stress \texttt{BAReferent\ +\ BAProposition}
against at least:

\begin{itemize}
\tightlist
\item
  behavior/event identity and specialized constraints in the
  facial-access branch;
\item
  information/resource lifecycle and correlation in order fulfillment;
\item
  internal/external responsibility boundaries plus contract/store
  candidates;
\item
  state/indeterminate semantics and physical handoff event reuse;
\item
  one bounded method-consumer projection.
\end{itemize}

Its purpose is to decide whether any deferred semantic kind is forced to
become a first-class split, whether the two-family candidate can be
simplified further, or whether it survives unchanged.

Do not start BA2 unless BA1 is explicitly closed after pressure testing
and closure review.

\end{document}
